\chapter{Conclusions}

The Midterm 1 milestone established the conceptual and analytical foundations of \textbf{Sample Browser}, resolving key product design questions through a Problem-Based Learning framework. By discarding our initial, unfeasible concept of an alternative consumer streaming shell, we focused on a genuine workflow bottleneck: independent producers losing hours scrubbing through long recordings to find isolated, sample-worthy passages.

Our domain research confirmed that contemporary streaming platforms index music strictly by whole-track metadata. This leaves creators trapped between generic, overused loop libraries like Splice and tedious, unstructured crate-digging across YouTube or Samplette. \textbf{Sample Browser} bridges this divide. By dividing tracks into 4--30 second passages, assigning structured acoustic tags, and using the YouTube IFrame API for compliant playback seeking and looping, our architecture protects the producer's creative focus.

This milestone yielded four tangible deliverables for our engineering roadmap:
\begin{itemize}
    \item \textbf{A validated problem definition} supported by academic MIR literature and empirical creator surveys (Tracklib, UK Musicians' Census).
    \item \textbf{A candid retrospective of project failures}, recording our team's pivot away from vocal minority feedback, passive user habits, and closed API barriers.
    \item \textbf{Four falsifiable validation hypotheses (H1--H4)} paired with a semi-structured interview protocol for field testing.
    \item \textbf{A competitive and SWOT benchmark} clearly positioning our application against Splice, Samplette, Tracklib, Sononym, and Cyanite.
\end{itemize}

With the domain analysis and system concept established, our team is prepared for Midterm 2. The next milestone will focus on validating our hypotheses through $N \approx 20$ recorded video interviews with active music creators (~4 per team member) backed by a supplementary quantitative survey, formalizing software requirements through UML modeling (Activity, Use Case, and Sequence Diagrams), and building interactive interface wireframes.
